% A sample source sheet showing the three Hebrew font roles. Needs Taamey Frank
% CLM, Taamey David CLM and Noto Rashi Hebrew installed; compile with XeLaTeX
% (twice) and commit the PDF, so readers see it without the fonts.
\documentclass[11pt]{article}
\usepackage[maggidshiur]{shiur}
\usepackage{fancyvrb}
\setshiurfont{mikra}{Taamey Frank CLM}
\setshiurfont{hebrew}{Taamey David CLM}
\setshiurfont{rashi}{Noto Rashi Hebrew}
\hypersetup{hidelinks}
\newfontfamily{\codefont}{FreeMono}
\fvset{formatcom=\codefont\small, frame=single, framesep=2mm}

\title{The light of the first day}
\author{A sample source sheet for the \texttt{shiur} package}
\date{}

\begin{document}

\makeTitle

\begin{note}
This sheet shows the three Hebrew roles, each in its own font: Tanakh (\texttt{mikra}) in
Taamey Frank CLM, which places trop well; rabbinic text (\texttt{hebrew}) in Taamey David
CLM; and a commentary in Rashi script (\texttt{rashi}) in Noto Rashi Hebrew. Notes like this
one appear only in the maggid shiur's version.
\end{note}

\begin{mekor}
\begin{mekorlang}{mikra}{בראשית א א-ג}
בְּרֵאשִׁ֖ית בָּרָ֣א אֱלֹהִ֑ים אֵ֥ת הַשָּׁמַ֖יִם וְאֵ֥ת הָאָֽרֶץ׃
וְהָאָ֗רֶץ הָיְתָ֥ה תֹ֙הוּ֙ וָבֹ֔הוּ וְחֹ֖שֶׁךְ עַל־פְּנֵ֣י תְה֑וֹם וְר֣וּחַ אֱלֹהִ֔ים מְרַחֶ֖פֶת עַל־פְּנֵ֥י הַמָּֽיִם׃
וַיֹּ֥אמֶר אֱלֹהִ֖ים יְהִ֣י א֑וֹר וַֽיְהִי־אֽוֹר׃
\end{mekorlang}
\begin{mekorlang}{english}{Genesis 1:1--3}
In the beginning God created the heavens and the earth. And the earth was without form and
void, and darkness was on the face of the deep, and the spirit of God hovered over the face of
the waters. And God said, ``Let there be light'', and there was light.
\end{mekorlang}
\end{mekor}

\begin{mekor}
\begin{mekorlang}{rashi}{רש״י על בראשית א א}
בראשית - אמר רבי יצחק: לא היה צריך להתחיל את התורה אלא מהחודש הזה לכם, שהיא מצוה ראשונה שנצטוו ישראל. ומה טעם פתח בבראשית? משום כח מעשיו הגיד לעמו לתת להם נחלת גוים.
\end{mekorlang}
\begin{mekorlang}{english}{Rashi on Genesis 1:1}
``In the beginning'' -- Rabbi Yitz\hdot ak said: the Torah need only have begun from ``This
month shall be for you'', the first commandment that Israel were given. Why, then, does it
begin with creation? Because ``He told His people the power of His works, to give them the
inheritance of the nations''.
\end{mekorlang}
\end{mekor}

\begin{note}
Why would the Torah need a reason to begin with creation at all?
\end{note}

\begin{mekor}
\begin{mekorlang}{hebrew}{בבלי חגיגה י״ב ע״א}
אמר רבי אלעזר: אור שברא הקדוש ברוך הוא ביום ראשון - אדם צופה בו מסוף העולם ועד סופו.
\end{mekorlang}
\begin{mekorlang}{english}{B. \Hdot agigah 12a}
Rabbi Elazar said: with the light that the Holy One, blessed be He, created on the first day,
a person could see from one end of the world to the other.
\end{mekorlang}
\end{mekor}

\section{How this sheet was made}

Each source names its role as its language, and the preamble chooses the fonts:

\begin{Verbatim}
\setshiurfont{mikra}{Taamey Frank CLM}
\setshiurfont{hebrew}{Taamey David CLM}
\setshiurfont{rashi}{Noto Rashi Hebrew}

\begin{mekorlang}{mikra}{בראשית א א-ג} ...
\begin{mekorlang}{rashi}{רש״י על בראשית א א} ...
\begin{mekorlang}{hebrew}{בבלי חגיגה י״ב ע״א} ...
\end{Verbatim}

\end{document}
